\documentclass[12pt]{article}

\usepackage[margin=1in]{geometry}
\usepackage{amsmath, amssymb}
\usepackage{graphicx}
\usepackage{booktabs}
\usepackage[hidelinks]{hyperref}
\usepackage[turkish]{babel}
\usepackage[titles]{tocloft}

\title{Yüksek Boyutlu Veri Setlerinde PCR (Principal Component Regression) Temel Bileşen Regresyonun Tahminsel Gücünün Değerlendirilmesi}
\author{Caner SELÇUK}
\date{\today}

\renewcommand{\cftsecleader}{\cftdotfill{\cftdotsep}}

\begin{document}
\shorthandoff{=}

\maketitle
\tableofcontents

\begin{abstract}
Bu raporda PCA'in, regresyon modellerinde bir uygulaması olan PCR'ı Beton Basınç Dayanımı \cite{dataset} veri seti üzerinde bir vaka çalışması olarak ele alınması ve çok boyutlu, çoklu doğrusal bağlantı (multicollinearity) içeren bağımsız değişkenleri daha düşük özniteliğe (feature) indirgeyip regresyon modellenmesi yapılmıştır. Ayrıca performans metrikleri incelenerek Çoklu Lineer Regresyon yöntemi ile karşılaştırılmıştır. Ridge/Lasso Regresyonları ile performans karşılaştırılması yapılmıştır. Son olarak yöntemin kısıtlılıkları tartışılıp alternatif yöntemlerden bahsedilmiştir.
\end{abstract}

\noindent\textbf{Anahtar Kelimeler:} PCR, Regresyon, Çoklu Doğrusal Bağlantı, Yüksek-Boyutlu Veri, PCA, Çoklu Doğrusal Regresyon, Öznitelik İndirgeme, Temel Bileşen Analiz, Temel Bileşen Regresyon

\newpage

\section{Giriş}
Günümüzde veri setlerinin öznitelik boyutları kolayca çok fazla boyutlara
ulaşabilmektedir. Örneğin; yüzlerce gen ifadesinin belirli tip bir kanserdeki ayırt etme etkisinin gözlemlenmesi, finansal karar verme süreçlerinde çok değiskenli veri setleri, kimyasal süreçlerde içerik miktarlarının optimal tahmin edilmesinde veya mobil uygulama yada yazılım hizmetlerinde kullanıcı deneyimi hakkında bilgi sahibi olmak için çok fazla parametrenin gözlemlenmesi. Buna genellikle ``Çok Boyutluluğun Laneti'' (Curse of Dimensionality) de denir. Çok boyutluluk beraberinde şu sorunları getirir:
veriyi görselliştirme zorluğu, modelin yorumlanabilirliğinin azalması, hedef değişkenin varyansını açıklamada etkisi olmayan bir yığın gereksiz değişkenler, çoklu doğrusal bağlantı ve mesafe tabanlı çalışan makine öğrenmesi algoritmaların (K En Yakın Komşu, Kümeleme vs.) düzgün çalışmamasına sebeb olur. Uzak-yakın kavramlarının daha az anlamı
olmasına sebebiyet verirler. Bağımsız değişkenler arasında güçlü korelasyona sahip değişkenlerin katsayıları tahmin edilirken belirsizliğin artmasına sebeb olur, katsayıların tahmin varyansını arttırırlar. Model veri setine aşırı uyum sağlar.

\vspace{1.5em}

Bu raporda ise PCR yöntemi ile $m \times n$ boyutlu tasarım matriksine PCA (Principal Component Analysis) uygulayarak boyut indirgemesi yapar. PCR denetimsiz bir modelleme yöntemdir çünkü hedef değişkeni ne kadar doğru tahmin ettiğini dikkate almaz. Boyut indirgemesi özniteliklerde en fazla varyansı yakalayan eksenleri bularak gerçekleşir. Tasarım matriksinde varyansı açıklayan en fazla $k$ bileşeni seçerek yeni bir regresyon modeli yapılır. Seçilecek $k$ hiper-parametresi çapraz doğrulama (cross-validation) yöntemiyle seçilir. Böylelikle çoklu doğrusal bağlantı içeren öznitelik sayısı yok edilir. Fakat yöntemin kısıtlı yanları adlı bölümde bu yöntemin yetersiz kaldığı durumları gösterip alternatif yöntemlerden bahsedildi. $F-test$ ve $R^2$ gibi istatistiksel metrikler ile indirgemenin hedef değişkeni tahmin etmede anlamlı fayda sağlayıp sağlamadığını test edildi. Modellemeler ve testler R programlama dili kullanılarak gerçekleştirildi.

\newpage

\section{Yöntem}
Describe the data, materials, and methods used in the study \cite{jolliffe}.

For example, an equation can be written as
\begin{equation}
  y = \alpha + \beta x + \varepsilon
\end{equation}

\section{Bulgular}
Present the main findings. Refer to tables and figures in the text.

\begin{table}[htbp]
    \centering
    \caption{Example results}
    \label{tab:results}
    \begin{tabular}{lcc}
        \toprule
        Measure & Group A & Group B \\
        \midrule
        Example measure & 1.23 & 4.56 \\
        \bottomrule
    \end{tabular}
\end{table}

\begin{figure}[htbp]
    \centering
    % Replace with your image file:
    \includegraphics[width=0.7\linewidth]{../plot}
    \caption{Example figure caption.}
    \label{fig:example}
\end{figure}

\section{Yöntemin Kısıtlı Yanları}
Interpret the findings, discuss their implications, and note limitations.

\section{Sonuç ve Öneriler}
Summarize the paper's main contribution and suggest directions for future
work.

\begin{thebibliography}{9}

    \bibitem{dataset}
    I-Cheng Yeh (1998).
    \textit{Concrete Compressive Strength [Dataset]}.
    UCI Machine Learning Repository, https://doi.org/10.24432/C5PK67.

    \bibitem{jolliffe}
    Jolliffe, I. T. (2002).
    \textit{Principal Component Analysis (2nd ed.)}.
    Springer-Verlag.

\end{thebibliography}

\end{document}

%%% Local Variables:
%%% mode: LaTeX
%%% TeX-master: t
%%% End:
